% C.41: root and FGH independent full proof/scope audits passed.
\subsection{The closed normal quotient by the descent-fixed subgroup}

\begin{statement}{C.41}{Proved}{Odd prime-power quotients after killing the fixed subgroup}
Put $G=\widehat{\mathrm{GT}}_0$ and let
$N_\varphi=\overline{\langle\!\langle C^\varphi\rangle\!\rangle}$
be the closed normal subgroup generated by the actual fixed subgroup.
There is a continuous surjection
\[
 G/N_\varphi\twoheadrightarrow
 C_2\times\prod_{\substack{p\geq5\ {
m prime}\\ k,j\geq1}}
                         \mathbb Z/p^k\mathbb Z.       \tag{C.41.1}
\]
For every $p\geq5$, its continuous $\mathbb F_p$ character
space has dimension exactly $\aleph_0$. In particular,
$C^\varphi$ has index $2^{\aleph_0}$ in $G$ and Haar
measure zero. No statement about the prime three part,
the complete abelianization, or a $\mathbb Z_p$ quotient
is implied. The inclusion $C^\varphi\subset\ker\Delta$
and the pentagon-to-fixedness question remain unresolved.
\end{statement}

\begin{proof}
We first make a locality point in C.38 explicit. Its
characters at a fixed prime $p$ are built on the actual
finite Magnus image of the local gentle group, together
with its local multiplier. The invariant homomorphism
of the multiplier kernel is extended across
$\mathbb Z_p^\times$, and the arithmetic correction
is a multiple of
\[
                  \log(\lambda_p^{p-1})/p\pmod{p^k}.
\]
Thus every character in that construction factors
through the single normalized local gentle group
$G_{0,p}$ of C.40. No other local coordinate or
global representative is used in its definition.

For fixed $p,k$, the construction gives arbitrarily
many jointly independent reductions on the actual
local multiplier kernel $U_p$. One may therefore
perform the final mixed-order selection of C.38
entirely inside these local characters: at each
step choose a character whose reduction on $U_p$
lies outside the span of all earlier reductions.
The relevant spaces have infinite dimension, so
this remains possible for every pair $(k,j)$.
For any finite collection, independence modulo $p$
gives $H+pT=T$ for its image $H$ in the finite
target $T$. Finite Nakayama gives $H=T$.
Compactness then gives actual continuous local
characters
\[
 \theta_{p,k,j}:G_{0,p}\longrightarrow\mathbb Z/p^k\mathbb Z
                                                        \tag{C.41.2}
\]
whose joint restriction to $U_p$ maps onto the
whole product over $k,j$. Their global inflations
kill the specified Galois image, as in C.38.

For each of these characters define on the full group
\[
 \eta_{p,k,j}(g)=
       \theta_{p,k,j}(g_p)-\theta_{p,k,j}(D(g)_p).
                                                        \tag{C.41.3}
\]
Both terms are continuous additive characters because
$D$ is a homomorphism. Equation (C.41.3) vanishes
pointwise on $C^\varphi=\operatorname{Fix}(D)$, hence
on its closed normal closure.

To check joint surjectivity, prescribe any finite
collection of target values. For each coefficient
prime involved, choose $\alpha_p\in U_p$ giving
those values under (C.41.2), and prescribe the
second local action to be the identity. C.40
supplies an actual full multiplier-one element
$g$ with $g_p=\alpha_p$ and $D(g)_p=1$ at all
these primes. Its values in (C.41.3) are precisely
the prescribed values. Compactness proves
surjectivity onto the entire odd-primary product
in (C.41.1).

There is also the quadratic character
\[
                       \eta_2=\Delta+\Delta\circ D.
                                                        \tag{C.41.4}
\]
It vanishes on the same fixed subgroup and its
closed normal closure. The actual element of
C.23 has $(\Delta(g),\Delta(D(g)))=(1,0)$,
so (C.41.4) is onto. For every finite target
subproduct, its quadratic factor and its odd
primary factors have coprime orders. Surjectivity
to each primary factor therefore gives
surjectivity to their product. A final compactness
argument proves (C.41.1).

The displayed quotient supplies infinitely many
independent mod-$p$ characters for every $p\geq5$.
The full group is countably based, as used in
C.38, so it has only countably many continuous
maps to a fixed finite group. This gives the
matching upper bound $\aleph_0$ for the dimension.

The product in (C.41.1) has cardinality
$2^{\aleph_0}$. Since $C^\varphi\subset N_\varphi$,
the coset space $G/C^\varphi$ maps onto this
quotient and has at least that cardinality.
The countably based compact group $G$ has at
most that cardinality. Its fixed subgroup
therefore has exactly the asserted index.
For Haar measure, arbitrarily large finite
quotients in (C.41.1) put $C^\varphi$ in
subgroups of arbitrarily large finite index.
Its measure is consequently zero.

This argument uses difference characters. It
does not identify $\Delta$ itself on the fixed
subgroup, and it does not say that the new
difference characters vanish on the full
pentagon group. Those implications would
require additional information about $D$ on
that subgroup.
\end{proof}
